\chapter{ОНОЛЫН БОЛОН ХАРЬЦУУЛСАН СУДАЛГАА}

\section{Хиймэл оюуны агентын ойлголт}

Хиймэл оюуны агент гэдэг нь орчноо хүлээн авч, түүн дээр үндэслэн шийдвэр гаргаж, үйлдэл хийх чадвартай систем юм \citep{russell-norvig}. Сүүлийн үеийн хэлний загваруудын хувьд энэ ойлголт нь шинэ утга агуулгатай болсон. Загвар нь зөвхөн текст үүсгээд зогсохгүй, гаднах системүүдтэй харилцах хэрэгслүүдийг дуудаж, олж авсан мэдээллийг нэгтгэн, цаашдын үйлдлийг төлөвлөх чадвартай болсон.

Орчин үеийн AI агент нь гурван үндсэн бүрэлдэхүүнээс бүрдэнэ. Нэгдүгээрт, хэлний загвар (LLM) нь хэрэглэгчийн хүсэлтийг ойлгох, хариулт үүсгэх, шийдвэр гаргах ``тархи'' болж үйлчилнэ. Хоёрдугаарт, хэрэгслүүд (tools) нь загварт гаднах системүүдтэй харилцах боломжийг олгодог ``гар, хөл'' юм. Гуравдугаарт, orchestration давхарга нь хэрэгслүүдийг зөв дарааллаар дуудах, үр дүнг нэгтгэх, алдаа засах үүрэгтэй.

\subsection{ReAct загвар}

AI агентын нэг түгээмэл архитектур нь ReAct (Reasoning and Acting) загвар юм \citep{react-paper}. Энэ загварын гол санаа нь загварыг ``бодолт'' болон ``үйлдэл''-ийг ээлжлэн хийхэд сургах явдал. Загвар нь эхлээд асуудлыг задлан шинжилж (reasoning), дараа нь шаардлагатай хэрэгслийг дуудаж (acting), хэрэгслийн хариуг авч дахин бодолт хийж, эцсийн хариулт бүрдүүлдэг.

Энэ загвар нь хэрэглэгчийн хүсэлтийг шууд нэг хэрэгсэлд холбохоос илүү уян хатан юм. Жишээлбэл, ``миний хамгийн сул дүнтэй хичээл юу вэ'' гэсэн асуулт нь эхлээд дүнгийн жагсаалтыг авах, дараа нь хамгийн багыг нь олох гэсэн хоёр алхамыг агуулж байна. ReAct загвар нь ийм олон алхамт асуултыг шат дараатай шийдвэрлэх боломжийг олгодог.

\subsection{Tool-calling механизм}

Орчин үеийн хэлний загварууд нь tool-calling буюу функц дуудах чадвартай болсон \citep{openai-functions}. Энэ нь загварт тодорхой схемтэй хэрэгслүүдийн жагсаалт өгөх, загвар хэрэглэгчийн хүсэлтэд тохирох хэрэгслийг сонгож дуудах, хэрэгслийн хариуг загварт буцааж өгч эцсийн хариулт үүсгэхийг хэлнэ.

Tool-calling нь AI системийг ``хаалттай'' байдлаас ``нээлттэй'' байдалд шилжүүлдэг. Загвар нь өөрийн сургалтын өгөгдөлд хязгаарлагдахгүй, бодит цагийн өгөгдөлд хандах, гаднах системүүдтэй харилцах боломжтой болдог. Гэхдээ энэ нь мөн аюулгүй байдлын шинэ асуудлуудыг бий болгодог тул read-only болон mutation хэрэгслүүдийг ялгах, баталгаажуулалтын механизм нэмэх зэрэг арга хэмжээ шаардлагатай.

\section{Model Context Protocol (MCP)}

Model Context Protocol нь Anthropic компаниас санаачилсан нээлттэй стандарт бөгөөд хэлний загварт гадаад хэрэгсэл, өгөгдлийн эх үүсвэрийг стандартчилсан аргаар холбох зориулалттай \citep{mcp-intro}. MCP-ийн гол зорилго нь интеграцын ``N x M'' асуудлыг шийдвэрлэх явдал юм. Өөрөөр хэлбэл, N тооны AI application болон M тооны гаднах системийг нэг бүрчлэн холбохын оронд, нийтлэг протокол ашиглан хоёр талаас нэг удаа хэрэгжүүлэхэд хангалттай болгох.

\subsection{MCP-ийн бүтэц}

MCP нь гурван үндсэн нэгжтэй. Tools нь загвараас дуудагдах функцууд бөгөөд тодорхой үйлдэл гүйцэтгэж үр дүн буцаадаг. Resources нь загварт өгөгдөл, контекст хангах файл, өгөгдлийн эх үүсвэрүүд юм. Prompts нь загварт заавар, загвар өгөх урьдчилан бэлтгэсэн текстүүд юм.

MCP сервер нь эдгээр нэгжүүдийг ил гаргадаг бие даасан процесс бөгөөд MCP client нь тэдгээрийг ашигладаг AI application юм. Энэ хоёр тал нь JSON-RPC протоколоор харилцана.

\subsection{Энэхүү төсөлд MCP сонгосон шалтгаан}

SISi Intelligent Agent төсөлд MCP сонгох хэд хэдэн шалтгаан байв. Нэгдүгээрт, стандарт интерфейс нь SISI, багшийн портал, календарь, баримт боловсруулах үйлчилгээ зэрэг өөр өөр системийг ижил загвараар илэрхийлэх боломжийг олгодог. Хоёрдугаарт, модульчлагдсан бүтэц нь порталтай холбогдох логик, Chrome extension-д суурилсан хэрэглэгчийн оролт, ирээдүйн интеграцуудыг тусдаа давхаргад хөгжүүлэх боломжтой болгодог. Гуравдугаарт, өргөтгөх боломж нь шинэ хэрэгсэл нэмэхэд бүх системийг дахин бичих шаардлагагүй болгодог. Дөрөвдүгээрт, аюулгүй байдлын удирдлага нь read-only болон mutation хэрэгслүүдийг тусад нь загварчлах боломжтой.

\section{Ижил төстэй системүүдийн судалгаа}

AI туслах системүүдийн өнөөгийн байдлыг гурван чиглэлээр авч үзэв: боловсролын AI туслахууд, tool-use AI системүүд, байгууллагын автоматжуулалтын шийдлүүд.

\subsection{Боловсролын AI туслахууд}

Боловсролын салбарт хиймэл оюуныг ашигласан хэд хэдэн систем байна. Duolingo-ийн AI туслах нь хэл сурахад харилцан ярилцах, алдаа засах, тайлбар өгөх чадвартай. Carnegie Learning-ийн MATHia платформ нь математикийн даалгавар шийдэхэд алхам алхмаар туслах, оюутны ахицыг хянах боломжтой \citep{carnegie-learning}. IBM Watson Education нь сургалтын материалыг хувь хүнд тохируулах, асуултад хариулах зэрэг функцтэй.

Эдгээр системүүдийн нийтлэг шинж нь сургалтын агуулгад төвлөрсөн байдал юм. Тэд хичээл заах, тайлбар өгөх, дасгал шалгах зэрэгт чадварлаг боловч сургуулийн дотоод порталын өгөгдөл, хэрэглэгчийн хувийн академик мэдээлэлтэй гүн интеграц хийх зорилгогүй. Энэ нь SISi Intelligent Agent-аас ялгаатай бөгөөд тус төсөл нь нийтлэг сургалтын агуулга бус, тухайн оюутны бодит өгөгдөлд суурилдаг.

\subsection{Tool-use AI системүүд}

Tool-calling чадвар бүхий AI системүүд сүүлийн жилүүдэд хурдацтай хөгжиж байна. OpenAI-ийн GPT Actions нь гуравдагч талын API-уудтай холбогдох боломжийг олгодог бөгөөд хэрэглэгч өөрийн хэрэгцээнд тохируулсан интеграц хийх боломжтой \citep{openai-actions}. Anthropic-ийн Claude нь MCP-ээр дамжуулан файл систем, өгөгдлийн сан, вэб хөтөч зэрэгтэй харилцах чадвартай болсон. GitHub Copilot нь код бичих, засахаас гадна terminal команд ажиллуулах, файл удирдах боломжтой болж байна.

Эдгээр системүүд нь ерөнхий зориулалттай бөгөөд олон төрлийн даалгаварт ашиглагдах боломжтой. Гэхдээ тэд тодорхой нэг домэйнд (жишээлбэл их сургуулийн портал) зориулсан биш тул тухайн домэйны нарийн шаардлагууд, workflow-уудыг бүрэн хангахгүй.

\subsection{Байгууллагын автоматжуулалтын системүүд}

Байгууллагын хэмжээнд ашиглагддаг AI туслах системүүд ч бас холбогдох судалгааны объект болно. Intercom Fin нь хэрэглэгчийн асуултад автоматаар хариулах, шаардлагатай үед хүн рүү шилжүүлэх чадвартай \citep{intercom-fin}. Zendesk AI нь support ticket-үүдийг ангилах, хариулт санал болгох, нийтлэг асуултад автоматаар хариулах боломжтой. ServiceNow Virtual Agent нь байгууллагын дотоод системүүдтэй холбогдож, ажилчдын хүсэлтийг автоматаар шийдвэрлэх чадвартай.

Эдгээр системүүд нь байгууллагын дотоод өгөгдөлтэй холбогдох чадвараараа SISi Intelligent Agent-тай төстэй. Гэхдээ тэд ихэвчлэн худалдааны бүтээгдэхүүн бөгөөд тухайн байгууллагын нарийн хэрэгцээнд тохируулахад хязгаарлалттай байдаг.

\subsection{Харьцуулалтын дүгнэлт}

\begin{table}[H]
  \centering
  \caption{Ижил төстэй системүүдийн харьцуулалт}
  \label{tab:similar-systems}
  \begin{tabularx}{\textwidth}{>{\raggedright\arraybackslash}p{3cm} >{\raggedright\arraybackslash}X >{\raggedright\arraybackslash}X}
    \toprule
    Системийн төрөл & Давуу тал & SISi Intelligent Agent-тай харьцуулахад \\
    \midrule
    Боловсролын AI (Duolingo, Carnegie Learning) & Сургалтын агуулгад чадварлаг, харилцан яриа сайн & Оюутны хувийн академик өгөгдөл, порталын интеграц байхгүй \\
    \addlinespace
    Tool-use AI (GPT Actions, Claude MCP) & Олон төрлийн системтэй холбогдох ерөнхий чадвар & Тодорхой домэйнд зориулсан биш, их сургуулийн portal workflow мэдэхгүй \\
    \addlinespace
    Байгууллагын AI (Intercom, ServiceNow) & Дотоод системтэй холбогдох, автоматжуулалт & Худалдааны бүтээгдэхүүн, тохируулах боломж хязгаарлагдмал \\
    \addlinespace
    SISi Intelligent Agent & NUM-ийн SISI порталд зориулсан, MCP-д суурилсан модульчлагдсан архитектур & Одоогоор хөгжүүлэлтийн шатанд, бүрэн туршилт хийгдээгүй \\
    \bottomrule
  \end{tabularx}
\end{table}

Энэхүү харьцуулалтаас харахад зах зээл дээрх системүүд тус бүрдээ давуу талтай боловч NUM-ийн SISI порталын нарийн шаардлагууд, оюутны өдөр тутмын workflow-уудад зориулсан систем байхгүй байна. Энэ нь SISi Intelligent Agent төслийн шинэлэг тал бөгөөд тодорхой домэйнд зориулсан AI агентыг MCP стандартаар загварчилсан нь бусад системээс ялгаатай онцлог болно.

\section{Судалгааны дүгнэлт}

Онолын болон харьцуулсан судалгааны үр дүнд хэд хэдэн дүгнэлтэд хүрэв.

Нэгдүгээрт, SISi Intelligent Agent төрлийн системд дан ганц LLM хангалтгүй бөгөөд бодит системтэй холбогдох агент архитектур шаардлагатай. Загварын оюун ухаан хангалттай боловч түүнд зохих хэрэгслүүдийг олгох шаардлагатай.

Хоёрдугаарт, Model Context Protocol нь тус төслийн хувьд интеграцын давхаргыг стандартчилах, өргөтгөх, аюулгүй удирдах хамгийн тохиромжтой шийдлийн нэг юм. MCP нь олон системийг нэг протоколоор холбох, read-only болон mutation хэрэгслүүдийг ялгах боломжийг олгодог.

Гуравдугаарт, ижил төстэй системүүд байдаг боловч NUM-ийн SISI порталын хэрэглээ, өгөгдөл, workflow-д чиглэсэн тусгай шийдэл хэрэгтэй. Ерөнхий зориулалтын tool-use AI системүүд нь энэ хэрэгцээг бүрэн хангахгүй.

Дөрөвдүгээрт, судалгааны дараагийн гол алхам нь архитектурын зөв сонголтыг бодит хэрэгжилт болон хэрэглэгчийн үнэлгээгээр баталгаажуулах явдал юм.
